\section{研究背景}
近年，情報通信技術の急速な発展に伴い，ネットワーク上のデータトラフィックは爆発的に増大している．
2024年版のEricsson Mobility Reportによると，世界のモバイルデータトラフィックは2024年時点で月間164エクサバイトに達しており，2031年に向けて年平均成長率約16\%で増加し続けると予測されている\cite{ericsson2024}．
特に5Gの普及により，2031年には全トラフィックの83\%が5G経由になると見込まれており，ネットワークインフラへの負荷は劇的に高まっている\cite{ericsson2024}．

また，通信の質に対する要求も高度化している．
メタバース，クラウドVR (Virtual Reality)，および自動運転といった次世代アプリケーションは，単なる接続性だけでなく，超広帯域かつ超低遅延な通信を必要とする．
例えば，没入感を維持するためのVRシステムでは，1Gbps以上の帯域幅に加え，動作から映像表示までを表す Motion-to-Photon レイテンシを20ms未満に抑えることが求められている\cite{sandvine2024}\cite{Yaqoob2020}．
さらに，メタバース空間におけるリアルタイムな相互作用を実現するためには，膨大な3次元データを遅延なく同期させる必要があり，ネットワークのQoS (Quality of Service) 保証がサービスの成否を分ける重要な要因となっている\cite{Hovan2021}．
加えて，IOWN (Innovative Optical and Wireless Network) 構想や，Beyond 5Gおよび6Gに向けた議論においては，電子処理を介さない All-Photonics Network (APN) による，エンドツーエンドでの遅延の極小化と帯域保証が重要なテーマとなっている\cite{iowngf2024}．

このような背景から，現代のネットワーク経路制御においては，従来のベストエフォート型ルーティングではなく，複数のQoSを同時に満たす高度な経路探索技術が不可欠となっている．

\section{ネットワーク経路制御の課題}
Open Shortest Path First (OSPF) や Border Gateway Protocol (BGP) などの従来の経路制御プロトコルでは，ホップ数や固定リンクコストといった単一の指標に基づいて最短経路を計算することが一般的であった\cite{Nurhaida2019}．
しかし，前述したような高精細映像伝送やリアルタイム制御を行う場合，以下の2つの要件を同時に満たす必要がある．

\begin{enumerate}
  \item ボトルネック帯域の最大化: 経路上で最も帯域が狭いボトルネックリンクの容量を最大化し，スループットを確保すること．
  \item 遅延制約の遵守: エンドツーエンドの遅延和が，アプリケーションの許容限界以下であること．
\end{enumerate}

このように，複数の制約条件を満たす経路を探索する問題は多制約経路問題 (MCOP:Multi-Constrained Optimal Path Problem) と呼ばれ，計算機科学においてNP困難な問題であることが知られている\cite{garey1979}．
特に，ボトルネック帯域のような最小値が全体を決める凹型メトリックと，遅延のような総和が全体を決める加法的メトリックという性質の異なる指標を同時に扱う場合，単純なダイクストラ法では最適解を導出できない．

厳密解法としてはラベル修正法などが提案されているが\cite{guerriero2001}，これらはパレート最適解を網羅的に探索するため，ネットワークのノード数や制約数が増加すると計算時間が指数関数的に増大するという課題がある．
ミリ秒単位の制御が求められる次世代ネットワークにおいて，計算に長時間を要する厳密解法を適用することは現実的ではない．

\section{研究目的}
上述の課題に対し，近年では計算コストの低いメタヒューリスティクスの適用が検討されている．
中でも粒子群最適化 (PSO) \cite{kennedy1995}は，実装が容易であり収束が速いことから，動的なネットワーク制御に適しているとされる．
先行研究において乗松らは，ボトルネックリンク最大化問題 (MBL:Maximum Bottleneck Link Problem) に対してPSOを適用し，その有効性を示した\cite{norimatsu2024}．

しかし，標準的なPSOには早熟収束という課題があり，複雑な制約条件下では局所最適解に停滞しやすく，真の最適解に到達できない場合がある．
また，実用化のためには，大規模なネットワークであっても数ミリ秒から数十ミリ秒程度で解を出力できる計算速度が求められる．

そこで本研究では，ボトルネックリンク帯域最大化を目的とした多制約経路問題に対し，以下の特徴を持つ改良型PSOを提案し，その有効性を評価することを目的とする．

\begin{enumerate}
  \item アルゴリズムの改良: 局所解回避能力の高いlBestトポロジーおよび探索停滞を検知して再探索を行うリスタート戦略を導入し，解の精度を向上させる．
  \item 並列化実装: マルチコアプロセッサを用いた並列処理により，厳密解法と比較して実用的な時間内での解探索を実現する．
\end{enumerate}

\section{本論文の構成}
本論文は全8章で構成される．
第2章では，本研究の基礎となるネットワークグラフ理論および最短経路問題(SPP:Shortest Path Problem)，ならびにMBL問題の定義について述べる．
第3章では，メタヒューリスティクスおよびPSOの基礎原理について解説する．
第4章では，関連する既存研究について概説し，本研究の位置付けを明確にする．
第5章では，提案手法であるlBestトポロジーとリスタート戦略を用いたPSOの詳細，およびその並列化実装について述べる．
第6章では，計算機シミュレーションによる評価実験の結果を示し，提案手法の有効性について考察する．
最後に第7章で本研究の結論を述べる．